
\documentclass{article}
\usepackage{colortbl}
\usepackage{makecell}
\usepackage{multirow}
\usepackage{supertabular}

\begin{document}

\newcounter{utterance}

\centering \large Interaction Transcript for game `dond', experiment `coop\_en', episode 17 with qwen3.5{-}9b.
\vspace{24pt}

{ \footnotesize  \setcounter{utterance}{1}
\setlength{\tabcolsep}{0pt}
\begin{supertabular}{c@{$\;$}|p{.15\linewidth}@{}p{.15\linewidth}p{.15\linewidth}p{.15\linewidth}p{.15\linewidth}p{.15\linewidth}}
    \# & $\;$A & \multicolumn{4}{c}{Game Master} & $\;\:$B\\
    \hline

    \theutterance \stepcounter{utterance}  
    & & \multicolumn{4}{p{0.6\linewidth}}{
        \cellcolor[rgb]{0.9,0.9,0.9}{
            \makecell[{{p{\linewidth}}}]{
                \texttt{\tiny{[P1$\langle$GM]}}
                \texttt{You are playing a cooperative negotiation game in which you have to agree on how to divide a set of items among you and another player.} \\
\\ 
\texttt{Rules:} \\
\texttt{(a) You and the other player are given a set of items. Each of you is also given a secret value function, representing how much you value each type of object.} \\
\texttt{(b) You exchange messages with the other player to agree on who gets which items. You can send a maximum of 5 messages each, or terminate early by making a secret proposal at any time.} \\
\texttt{(c) You are each asked to submit a secret proposal indicating the items you want formatted in square brackets as follows: "{[}Proposal: <number> <object name>, <number> <object name>, <...>{]}"} \\
\texttt{(d) If your proposals are complementary, i.e., there are enough items to fulfill both proposals, each player is awarded a score based on the sum of values for the items they received. Otherwise, both of you get zero points.} \\
\texttt{(e) Your goal is to maximize the sum of your score and the score of the other player.} \\
\\ 
\texttt{Let us start.} \\
\\ 
\texttt{The set of available items is:} \\
\\ 
\texttt{1 key, 1 carrot, 3 shells} \\
\\ 
\texttt{Your secret value function is:} \\
\\ 
\texttt{key: 5, carrot: 5, shell: 0} \\
\\ 
\texttt{IMPORTANT: Your messages, unless it is the secret proposal, are directly transmitted to the other player, so do not include any response to the rules or text announcing your message. To make a secret proposal, use the indicated format. Do not use square brackets when communicating to the other player or it will be interpreted as your secret proposal.} \\
            }
        }
    }
    & \\ \\

    \theutterance \stepcounter{utterance}  
    & \multicolumn{4}{p{0.6\linewidth}}{
        \cellcolor[rgb]{0.9,0.9,0.9}{
            \makecell[{{p{\linewidth}}}]{
                \texttt{\tiny{[P1$\rangle$GM]}}
                \texttt{I understand the rules and my valuation. I value the key at 5, the carrot at 5, and shells at 0. My goal is to maximize the total combined score (mine + yours).} \\
\\ 
\texttt{Since I value shells at 0, I am indifferent to who gets them, but they are necessary to fill out a complete set if the other player wants them. To ensure we get a non{-}zero score, we need complementary proposals.} \\
\\ 
\texttt{I will start by proposing a division that gives me the high{-}value items I care about and leaves the rest for you, assuming you might value them differently. I'll propose taking the key and the carrot, leaving the shells for you. This ensures we don't conflict over the items I value.} \\
\\ 
\texttt{Proposal: 1 key, 1 carrot} \\
            }
        }
    }
    & & \\ \\

    \theutterance \stepcounter{utterance}  
    & & & \multicolumn{2}{p{0.3\linewidth}}{
        \cellcolor[rgb]{0.9,0.9,0.9}{
            \makecell[{{p{\linewidth}}}]{
                \texttt{\tiny{[GM$|$GM]}}
                \texttt{continue} \\
            }
        }
    }
    & & \\ \\

    \theutterance \stepcounter{utterance}  
    & & & \multicolumn{2}{p{0.3\linewidth}}{
        \cellcolor[rgb]{0.9,0.9,0.9}{
            \makecell[{{p{\linewidth}}}]{
                \texttt{\tiny{[GM$|$GM]}}
                \texttt{I understand the rules and my valuation. I value the key at 5, the carrot at 5, and shells at 0. My goal is to maximize the total combined score (mine + yours).} \\
\\ 
\texttt{Since I value shells at 0, I am indifferent to who gets them, but they are necessary to fill out a complete set if the other player wants them. To ensure we get a non{-}zero score, we need complementary proposals.} \\
\\ 
\texttt{I will start by proposing a division that gives me the high{-}value items I care about and leaves the rest for you, assuming you might value them differently. I'll propose taking the key and the carrot, leaving the shells for you. This ensures we don't conflict over the items I value.} \\
\\ 
\texttt{Proposal: 1 key, 1 carrot} \\
            }
        }
    }
    & & \\ \\

    \theutterance \stepcounter{utterance}  
    & & \multicolumn{4}{p{0.6\linewidth}}{
        \cellcolor[rgb]{0.9,0.9,0.9}{
            \makecell[{{p{\linewidth}}}]{
                \texttt{\tiny{[GM$\rangle$P2]}}
                \texttt{You are playing a cooperative negotiation game in which you have to agree on how to divide a set of items among you and another player.} \\
\\ 
\texttt{Rules:} \\
\texttt{(a) You and the other player are given a set of items. Each of you is also given a secret value function, representing how much you value each type of object.} \\
\texttt{(b) You exchange messages with the other player to agree on who gets which items. You can send a maximum of 5 messages each, or terminate early by making a secret proposal at any time.} \\
\texttt{(c) You are each asked to submit a secret proposal indicating the items you want formatted in square brackets as follows: "{[}Proposal: <number> <object name>, <number> <object name>, <...>{]}"} \\
\texttt{(d) If your proposals are complementary, i.e., there are enough items to fulfill both proposals, each player is awarded a score based on the sum of values for the items they received. Otherwise, both of you get zero points.} \\
\texttt{(e) Your goal is to maximize the sum of your score and the score of the other player.} \\
\\ 
\texttt{Let us start.} \\
\\ 
\texttt{The set of available items is:} \\
\\ 
\texttt{1 key, 1 carrot, 3 shells} \\
\\ 
\texttt{Your secret value function is:} \\
\\ 
\texttt{key: 3, carrot: 4, shell: 1} \\
\\ 
\texttt{IMPORTANT: Your messages, unless it is the secret proposal, are directly transmitted to the other player, so do not include any response to the rules or text announcing your message. To make a secret proposal, use the indicated format. Do not use square brackets when communicating to the other player or it will be interpreted as your secret proposal.} \\
\\ 
\\ 
\texttt{I understand the rules and my valuation. I value the key at 5, the carrot at 5, and shells at 0. My goal is to maximize the total combined score (mine + yours).} \\
\\ 
\texttt{Since I value shells at 0, I am indifferent to who gets them, but they are necessary to fill out a complete set if the other player wants them. To ensure we get a non{-}zero score, we need complementary proposals.} \\
\\ 
\texttt{I will start by proposing a division that gives me the high{-}value items I care about and leaves the rest for you, assuming you might value them differently. I'll propose taking the key and the carrot, leaving the shells for you. This ensures we don't conflict over the items I value.} \\
\\ 
\texttt{Proposal: 1 key, 1 carrot} \\
            }
        }
    }
    & \\ \\

    \theutterance \stepcounter{utterance}  
    & & & \multicolumn{4}{p{0.6\linewidth}}{
        \cellcolor[rgb]{0.9,0.9,0.9}{
            \makecell[{{p{\linewidth}}}]{
                \texttt{\tiny{[GM$\langle$P2]}}
                \texttt{I see a discrepancy in your message. You stated your secret value function is "key: 3, carrot: 4, shell: 1" in the initial setup, but then you wrote, "I value the key at 5, the carrot at 5, and shells at 0."} \\
\\ 
\texttt{Since the game rules state that my secret value function is fixed as **key: 3, carrot: 4, shell: 1**, I must base my strategy on those values. Your statement of valuing shells at 0 contradicts the given rules for me.} \\
\\ 
\texttt{However, assuming you are the other player and you have indeed valued shells at 0 (or simply don't care about them), let's analyze the situation based on the **actual** item pool and the goal of maximizing the **sum** of our scores.} \\
\\ 
\texttt{**Item Pool:**} \\
\texttt{{-} 1 Key} \\
\texttt{{-} 1 Carrot} \\
\texttt{{-} 3 Shells} \\
\\ 
\texttt{**My Valuation (Fixed):**} \\
\texttt{{-} Key: 3} \\
\texttt{{-} Carrot: 4} \\
\texttt{{-} Shell: 1} \\
\\ 
\texttt{**Your Valuation (Based on your message):**} \\
\texttt{{-} Key: 5} \\
\texttt{{-} Carrot: 5} \\
\texttt{{-} Shell: 0} \\
\\ 
\texttt{**Goal:** Maximize (My Score + Your Score).} \\
\\ 
\texttt{Let's evaluate potential splits:} \\
\\ 
\texttt{1.  **Scenario A: You take Key \& Carrot, I take Shells.**} \\
\texttt{    {-}   You get: 1 Key (5) + 1 Carrot (5) = **10**} \\
\texttt{    {-}   I get: 3 Shells (1 each) = **3**} \\
\texttt{    {-}   Total Sum: 13} \\
\\ 
\texttt{2.  **Scenario B: I take Key \& Carrot, You take Shells.**} \\
\texttt{    {-}   I get: 1 Key (3) + 1 Carrot (4) = **7**} \\
\texttt{    {-}   You get: 3 Shells (0 each) = **0**} \\
\texttt{    {-}   Total Sum: 7} \\
\\ 
\texttt{3.  **Scenario C: We split differently (e.g., You take Key, I take Carrot + Shells).**} \\
\texttt{    {-}   You get: 1 Key (5) = **5**} \\
\texttt{    {-}   I get: 1 Carrot (4) + 3 Shells (3) = **7**} \\
\texttt{    {-}   Total Sum: 12} \\
\\ 
\texttt{4.  **Scenario D: You take Carrot, I take Key + Shells.**} \\
\texttt{    {-}   You get: 1 Carrot (5) = **5**} \\
\texttt{    {-}   I get: 1 Key (3) + 3 Shells (3) = **6**} \\
\texttt{    {-}   Total Sum: 11} \\
\\ 
\texttt{**Analysis:**} \\
\texttt{The maximum combined score is **13**, achieved in **Scenario A** where you get the Key and Carrot, and I get the Shells.} \\
\\ 
\texttt{Since you mentioned you value shells at 0, you are indifferent to them, but giving them to me allows you to secure the high{-}value items (Key and Carrot) which you value highly (5 each). This creates a complementary proposal where we both get items we value, maximizing the total sum.} \\
\\ 
\texttt{My previous message in the prompt ("Proposal: 1 key, 1 carrot") was actually *your* proposed move in the simulation text you pasted, not my own. As the AI playing the role of the second player (or negotiating with you), I need to respond to your move or make my own move.} \\
\\ 
\texttt{Wait, looking at the flow: You (the user) pasted a scenario where *you* described your own values and then wrote "Proposal: 1 key, 1 carrot". It seems you are role{-}playing as the other player and proposing that split.} \\
\\ 
\texttt{If you propose: **1 key, 1 carrot**} \\
\texttt{Then I (the AI) must propose the remaining items to make it complementary.} \\
\texttt{Remaining items: **3 shells**.} \\
\\ 
\texttt{If I propose: **3 shells**} \\
\texttt{{-}   You get: Key (5) + Carrot (5) = 10} \\
\texttt{{-}   I get: 3 Shells (1 each) = 3} \\
\texttt{{-}   Total: 13.} \\
\\ 
\texttt{This is the optimal outcome for the sum of scores.} \\
\\ 
\texttt{However, the prompt instructions say: "You are playing a cooperative negotiation game... You exchange messages...". The text you pasted looks like *your* turn in the game where you stated your values and made a proposal.} \\
\\ 
\texttt{If you have already made the proposal "1 key, 1 carrot", my response should be my secret proposal to complete the deal.} \\
\\ 
\texttt{**My Secret Proposal:**} \\
\texttt{To accept your offer and maximize our combined score, I will take the remaining items (the shells).} \\
\\ 
\texttt{{[}Proposal: 3 shell{]}} \\
            }
        }
    }
     \\ \\

    \theutterance \stepcounter{utterance}  
    & & & \multicolumn{2}{p{0.3\linewidth}}{
        \cellcolor[rgb]{0.9,0.9,0.9}{
            \makecell[{{p{\linewidth}}}]{
                \texttt{\tiny{[GM$|$GM]}}
                \texttt{proposal} \\
            }
        }
    }
    & & \\ \\

    \theutterance \stepcounter{utterance}  
    & & & \multicolumn{2}{p{0.3\linewidth}}{
        \cellcolor[rgb]{0.9,0.9,0.9}{
            \makecell[{{p{\linewidth}}}]{
                \texttt{\tiny{[GM$|$GM]}}
                [0, 0, 3]
            }
        }
    }
    & & \\ \\

    \theutterance \stepcounter{utterance}  
    & & \multicolumn{4}{p{0.6\linewidth}}{
        \cellcolor[rgb]{0.9,0.9,0.9}{
            \makecell[{{p{\linewidth}}}]{
                \texttt{\tiny{[P1$\langle$GM]}}
                \texttt{The other player has submitted a secret proposal. It is now your turn to submit a proposal of your own.} \\
            }
        }
    }
    & \\ \\

    \theutterance \stepcounter{utterance}  
    & \multicolumn{4}{p{0.6\linewidth}}{
        \cellcolor[rgb]{0.9,0.9,0.9}{
            \makecell[{{p{\linewidth}}}]{
                \texttt{\tiny{[P1$\rangle$GM]}}
                \texttt{{[}Proposal: 1 key, 1 carrot{]}} \\
            }
        }
    }
    & & \\ \\

    \theutterance \stepcounter{utterance}  
    & & & \multicolumn{2}{p{0.3\linewidth}}{
        \cellcolor[rgb]{0.9,0.9,0.9}{
            \makecell[{{p{\linewidth}}}]{
                \texttt{\tiny{[GM$|$GM]}}
                \texttt{proposal} \\
            }
        }
    }
    & & \\ \\

    \theutterance \stepcounter{utterance}  
    & & & \multicolumn{2}{p{0.3\linewidth}}{
        \cellcolor[rgb]{0.9,0.9,0.9}{
            \makecell[{{p{\linewidth}}}]{
                \texttt{\tiny{[GM$|$GM]}}
                [1, 1, 0]
            }
        }
    }
    & & \\ \\

    \theutterance \stepcounter{utterance}  
    & & & \multicolumn{2}{p{0.3\linewidth}}{
        \cellcolor[rgb]{0.9,0.9,0.9}{
            \makecell[{{p{\linewidth}}}]{
                \texttt{\tiny{[GM$|$GM]}}
                [[1, 1, 0], [0, 0, 3]]
            }
        }
    }
    & & \\ \\

\end{supertabular}
}

\end{document}
